\chapter{Five-Reviewer Editorial Audit, Rebuttal, and Revision Traceability}
\label{app:reviewer-gap-analysis}

This appendix records a simulated five-reviewer R1-style evaluation of the ORIUS monograph. The point is not artificial volume; it is to expose the manuscript to the five lenses most likely to challenge a universal physical-AI safety thesis: controls, uncertainty, formal verification, systems runtime engineering, and cross-domain thesis coherence. The tables in this appendix are editorial audit aids for revision traceability; they are not part of the scientific evidence surface.

\begin{longtable}{p{0.22\textwidth}p{0.10\textwidth}p{0.58\textwidth}}
\caption{Near-final reviewer scorecard summary for the ORIUS monograph.}\label{tab:reviewer-scorecards}\\
\toprule
\textbf{Reviewer} & \textbf{Score} & \textbf{Primary concern}\\
\midrule
\endfirsthead
\toprule
\textbf{Reviewer} & \textbf{Score} & \textbf{Primary concern}\\
\midrule
\endhead
Formal Safety / Control Theory & A- & Equal-domain rhetoric must remain subordinate to the governed parity gate.\\
ML Uncertainty / Calibration & B+ & Coverage and calibration language must stay bounded by the actual replay evidence.\\
CPS / Systems and Runtime Architecture & A & Runtime and governance claims must remain concrete and artifact-backed across domains.\\
Physical-AI Deployment and Domain Safety & B+ & Navigation and aerospace still block any equal-peer deployment reading.\\
R1 Dissertation Committee Reader & A & The book is structurally strong, but evidence asymmetry must remain explicit everywhere.\\
\bottomrule
\end{longtable}

\section{Reviewer 1: Controls and Runtime Assurance}
\textbf{Summary judgment}: The monograph is scientifically credible as a universal runtime-safety architecture, but equal-domain universality remains gated because navigation and aerospace do not yet clear the same defended closure surface as the stronger rows.

\textbf{Scorecard verdict}: Strong R1 thesis; equal-domain flagship claim still gated by navigation and aerospace.

\textbf{Strongest contributions}
\begin{itemize}
\item The degraded-observation hazard is identified as a plant-agnostic safety object rather than a battery-specific modeling issue.
\item The theorem bridge is explicit about what is structural, what is bounded, and what must still be validated by replay or artifact.
\item The parity matrix keeps the strongest claim subordinate to governed evidence rather than rhetorical flattening.
\end{itemize}
\textbf{Major weaknesses}
\begin{itemize}
\item Battery still carries the deepest theorem-to-artifact lineage, so the equal-domain narrative must remain carefully bounded.
\item The formal surface remains much stronger for some domains than others.
\end{itemize}
\textbf{Unsupported claims}
\begin{itemize}
\item Any statement implying equal theorem-grade closure across all six domains.
\item Any statement implying universal hardware or regulatory readiness.
\end{itemize}
\textbf{Required experiments}
\begin{itemize}
\item Real-data navigation closure under the universal replay contract.
\item A stronger aerospace task with material post-repair improvement on the governing safety object.
\end{itemize}
\textbf{Required writing changes}
\begin{itemize}
\item Keep the theorem/evidence boundary stable in the theory bridge, synthesis, and conclusion chapters.
\item State the parity gate once and reuse it consistently.
\end{itemize}
\textbf{Decision recommendation}: Weak accept for thesis submission; reject any equal-domain flagship claim until the gated rows are closed.

\section{Reviewer 2: Statistical ML and Conformal Inference}
\textbf{Summary judgment}: The uncertainty layer is valuable because it treats reliability-conditioned uncertainty as a runtime safety object, but the book must remain precise about where coverage language is formal and where it is empirical or conservative.

\textbf{Scorecard verdict}: The uncertainty layer is credible; the monograph must still stop short of equal-domain statistical closure.

\textbf{Strongest contributions}
\begin{itemize}
\item The monograph ties conformal and calibration ideas to runtime action legality rather than to prediction quality alone.
\item The universal benchmark schema makes true-state violation and observed-state satisfaction separable across domains.
\item The non-i.i.d. and distribution-shift caveats are visible instead of hidden.
\end{itemize}
\textbf{Major weaknesses}
\begin{itemize}
\item The weaker rows do not yet support equal-domain statistical closure.
\item Coverage-style language can still be overread if chapter summaries are not carefully bounded.
\end{itemize}
\textbf{Unsupported claims}
\begin{itemize}
\item Universal statistical calibration under arbitrary shift.
\item Formal coverage guarantees for domains whose replay surface remains portability-only or experimental.
\end{itemize}
\textbf{Required experiments}
\begin{itemize}
\item Fault-mode and domain-specific coverage diagnostics for every defended row.
\item Additional uncertainty summaries for navigation and aerospace once their evidence surfaces are upgraded.
\end{itemize}
\textbf{Required writing changes}
\begin{itemize}
\item Differentiate bounded calibration claims from conservative widening heuristics.
\item Keep uncertainty claims tied to tracked replay artifacts.
\end{itemize}
\textbf{Decision recommendation}: Accept for a bounded universal-safety monograph; major revision for any stronger cross-domain statistical claim.

\section{Reviewer 3: Systems and Runtime Engineering}
\textbf{Summary judgment}: The systems contribution is strong because runtime behavior, fallback, certificates, and audit continuity are first-class objects in the main manuscript rather than hidden implementation details.

\textbf{Scorecard verdict}: Systems contribution is strong enough for serious review once the parity matrix stays central.

\textbf{Strongest contributions}
\begin{itemize}
\item CertOS is integrated into the scientific argument rather than treated as release tooling.
\item The replay and benchmark contract is specific enough to support cross-domain inspection.
\item The backend artifact path and governed publication surface keep the evidence path auditable.
\end{itemize}
\textbf{Major weaknesses}
\begin{itemize}
\item Runtime and fallback evidence are still deepest in the witness and strongest defended rows.
\item Cross-domain composition and lifecycle coverage remain selective.
\end{itemize}
\textbf{Unsupported claims}
\begin{itemize}
\item Full deployment readiness in every domain.
\item Uniform multi-agent or shared-constraint closure across all rows.
\end{itemize}
\textbf{Required experiments}
\begin{itemize}
\item Broader runtime-budget and lifecycle traces for the weaker domains.
\item Expanded cross-domain composition scenarios where the adapter contract actually supports them.
\end{itemize}
\textbf{Required writing changes}
\begin{itemize}
\item Keep latency, fallback, certificate validity, and audit continuity in the mainline claims.
\item Avoid relegating critical runtime semantics to appendices only.
\end{itemize}
\textbf{Decision recommendation}: Strong accept as a systems-and-governance thesis; bounded revision required before broader deployment claims.

\section{Reviewer 4: Cross-Domain Physical AI and Deployment}
\textbf{Summary judgment}: ORIUS has credible societal potential as a fundamental safety layer for physical AI, but that potential only reads as serious when the monograph keeps weaker domains explicitly gated instead of rhetorically promoted.

\textbf{Scorecard verdict}: The physical-AI safety layer claim is compelling when bounded deployment language is preserved.

\textbf{Strongest contributions}
\begin{itemize}
\item The book treats battery, AV, industrial, healthcare, navigation, and aerospace as first-class domain chapters rather than as add-on examples.
\item The parity gate prevents the deployment narrative from outrunning the artifact surface.
\item The domain chapters share one template, which makes the universality claim inspectable.
\end{itemize}
\textbf{Major weaknesses}
\begin{itemize}
\item Navigation and aerospace remain below equal-peer maturity.
\item Some readers will still look for broader field closure than the current evidence supports.
\end{itemize}
\textbf{Unsupported claims}
\begin{itemize}
\item Equal defended deployment credibility across every domain.
\item Aerospace or navigation as mature peers to the witness and strongest defended rows.
\end{itemize}
\textbf{Required experiments}
\begin{itemize}
\item A defended real-data navigation row.
\item A non-placeholder aerospace benchmark with a stronger safety object and post-repair gain.
\end{itemize}
\textbf{Required writing changes}
\begin{itemize}
\item Retain explicit domain non-claims in every chapter.
\item Keep deployment language bounded by the parity matrix and artifact register.
\end{itemize}
\textbf{Decision recommendation}: Accept as a universal physical-AI safety architecture; major revision for any flat equal-domain deployment claim.

\section{Reviewer 5: R1 Dissertation and Monograph Coherence}
\textbf{Summary judgment}: The manuscript now reads like a real R1-style monograph because it has one hazard, one architecture, one theorem bridge, six domain chapters, and one parity gate. The remaining gap is evidence symmetry, not book structure.

\textbf{Scorecard verdict}: Book quality is near submission-ready; the remaining gap is evidence parity, not structure.

\textbf{Strongest contributions}
\begin{itemize}
\item The main body no longer depends on stitched paper-lineage framing.
\item The universal-first narrative is visible from the abstract through the conclusion.
\item The appendices provide bibliography, review, protocol, crosswalk, and artifact traceability at book scale.
\end{itemize}
\textbf{Major weaknesses}
\begin{itemize}
\item Legacy repo surfaces can still contradict the monograph if they are not treated as archival support only.
\item The equal-domain ambition is stronger than the current parity evidence.
\end{itemize}
\textbf{Unsupported claims}
\begin{itemize}
\item A claim that the monograph has already earned equal-domain closure.
\item Any suggestion that weaker rows are present only for symmetry rather than bounded scientific use.
\end{itemize}
\textbf{Required experiments}
\begin{itemize}
\item Only the parity-closing experiments needed for navigation and aerospace if equal-domain promotion remains the target.
\end{itemize}
\textbf{Required writing changes}
\begin{itemize}
\item Keep the main body universal-first and the weaker rows honest.
\item Preserve the explicit distinction between architecture universality and evidence parity in the introduction, synthesis, and conclusion.
\end{itemize}
\textbf{Decision recommendation}: Accept as a thesis-length monograph with bounded universal claims; defer stronger flagship universality language until parity is actually earned.

\section{Consolidated Editorial Findings}
\begin{longtable}{p{0.22\textwidth}p{0.24\textwidth}p{0.18\textwidth}p{0.26\textwidth}}
\caption{Consolidated editorial findings across the five reviewer lenses.}\label{tab:reviewer-gap-matrix}\\
\toprule
\textbf{Gap} & \textbf{Where it appears} & \textbf{Reviewer pressure} & \textbf{Revision action}\\
\midrule
\endfirsthead
\toprule
\textbf{Gap} & \textbf{Where it appears} & \textbf{Reviewer pressure} & \textbf{Revision action}\\
\midrule
\endhead
Evidence asymmetry across domains & Parity matrix, cross-domain synthesis, and domain status language & Controls, deployment, and committee readers all insist that architecture universality must not be confused with equal-domain closure & Keep battery as witness, AV plus industrial and healthcare as defended bounded rows, navigation as shadow-synthetic, and aerospace as experimental until stronger reruns exist.\\
Coverage and calibration ambiguity & Related work, uncertainty discussion, and benchmark interpretation & Statistical-ML review demands a clear line between formal calibration claims and conservative widening heuristics & State where coverage language is theorem-backed, where it is replay-backed, and where it is only heuristic.\\
Theorem / replay / implementation drift & Theory bridge, synthesis, and conclusion & Formal review demands an explicit claim boundary that is reused verbatim & Tie every strong claim to a theorem object, tracked artifact, or explicit bounded empirical result.\\
Runtime and audit visibility & Governance, latency, and fallback chapters & Systems review requires runtime budgets, certificate lifecycle, and audit continuity to stay in the main narrative & Keep governance and fallback semantics first-class in the architecture and synthesis chapters.\\
Legacy framing leakage & Non-canonical repo docs and archival artifacts & Committee and deployment readers will downgrade the book if older battery-origin or stitched-paper language leaks into the main story & Keep the canonical monograph universal-first and treat archival paper-lineage artifacts as support only.\\
\bottomrule
\end{longtable}

\section{Author Rebuttal}
ORIUS does not need to claim equal empirical maturity in every domain to justify the monograph. It needs to justify one universal runtime safety layer and then state, with governed discipline, how each domain currently maps onto that contract.
\begin{itemize}
\item \textbf{To the controls reviewer}: the monograph keeps the parity gate reader-visible and does not blur architecture with equal closure.
\item \textbf{To the statistical-ML reviewer}: calibration language is bounded to the evidence surface and never upgraded into universal statistical guarantees by prose.
\item \textbf{To the systems reviewer}: runtime behavior, fallback, certificate validity, and audit continuity remain part of the main scientific contribution rather than appendix-only material.
\item \textbf{To the deployment reviewer}: weaker rows remain explicitly gated so the physical-AI safety-layer claim does not borrow credibility from unclosed domains.
\item \textbf{To the committee reviewer}: the book keeps one universal narrative spine from hazard to parity gate and uses the appendices only to deepen, not replace, the main argument.
\end{itemize}

\section{Revision Traceability}
\begin{longtable}{p{0.16\textwidth}p{0.34\textwidth}p{0.30\textwidth}}
\caption{Revision traceability matrix from reviewer critique to manuscript action.}\label{tab:reviewer-traceability}\\
\toprule
\textbf{Reviewer} & \textbf{Critique} & \textbf{Revision action / file target}\\
\midrule
\endfirsthead
\toprule
\textbf{Reviewer} & \textbf{Critique} & \textbf{Revision action / file target}\\
\midrule
\endhead
Controls and runtime assurance & Do not blur universal architecture with equal-domain parity. & Keep the parity gate central in \texttt{monograph/ch14\_cross\_domain\_synthesis.tex} and align all domain labels to the tracked parity matrix.\\
Statistical ML / conformal & Separate formal calibration from conservative widening. & Keep claim-boundary language explicit in \texttt{monograph/ch02\_oasg\_claim\_boundary.tex} and the term register.\\
Systems and runtime engineering & Keep governance, fallback, latency, and certificate lifecycle visible. & Preserve runtime and governance emphasis in \texttt{monograph/ch07\_system\_benchmark\_governance.tex} and the supporting CertOS chapters.\\
Cross-domain deployment & Keep navigation and aerospace gated until stronger evidence exists. & Maintain explicit non-claims and promotion obligations in the domain chapters and protocol-card appendix.\\
R1 monograph coherence & Preserve one universal-first narrative spine and treat archival language as non-canonical. & Keep \texttt{paper/paper.tex}, the monograph chapters, and the authoring guide synchronized around the same universal-first structure.\\
\bottomrule
\end{longtable}

